% !PRINTED 105
\weberpage{105}
%<<CONT
und erhalten, wenn $\varepsilon_1$, $\varepsilon_2$ den Werth 0 oder 1 haben, die Resolventen

\weq{6}{\begin{aligned}
\Omega_{\varepsilon_1,\varepsilon_2} &= \sum^n (-1)^{\varepsilon_1\sigma_1+\varepsilon_2\sigma_2}\, r^n\\
&= \eta + (-1)^{\varepsilon_1}\eta' + (-1)^{\varepsilon_2}\eta'' + (-1)^{\varepsilon_1+\varepsilon_2}\eta'''.
\end{aligned}}

Ist nun $m = q_1\, q_2 \ldots q_\mu$ nicht durch 8 theilbar, so können
wir, wenn $r_1, r_2, \ldots r_\mu$ primitive Wurzeln der Grade $q_1, q_2, \ldots q_\mu$
bedeuten,
\[ r = r_1\, r_2 \ldots r_\mu \]
setzen, ferner
\[ \begin{aligned}
n &\equiv n_1\ (\mathrm{mod}\ q_1)\\
&\equiv n_2\ (\mathrm{mod}\ q_2)\\
&\cdot\ \cdot\ \cdot\ \cdot\ \cdot\ \cdot\ \cdot\\
&\equiv n_\mu\ (\mathrm{mod}\ q_\mu),
\end{aligned} \]
wodurch sich ergiebt:
\weq{7}{\Omega_{\varepsilon_1,\varepsilon_2} = \prod_{1,\mu}^{h} \sum^{n_h} (-1)^{\nu_h(\varepsilon_1 e_{1,h}+\varepsilon_2 e_{2,h})}\, r_h^{n_h}.}

Die einzelnen Factoren dieses Productes sind die in \S.\,171
(Bd.\ I) bestimmten \gesperrt{Gauss'schen Summen}, die dort durch
Quadratwurzeln ausgedrückt sind. Nur wenn $q_1 = 4$ ist und
$\varepsilon_1\, e_{1,1} + \varepsilon_2\, e_{1,2} = 0$, ist $\Omega_{\varepsilon_1,\varepsilon_2} = 0$. % [sic: Weber prints e_{1,2}; specializing the preceding general factor to h=1 requires e_{2,1}.]

Ist aber $m$ durch 8 theilbar und
\[ m = q_2\, q_3 \ldots q_\mu, \quad q_2 = 8, \]
so erhält man, wenn man $r = r_2\, r_3 \ldots r_\mu$ setzt:
\[ \begin{aligned}
\Omega_{\varepsilon_1,\varepsilon_2} &= \sum^{n_2} (-1)^{\varepsilon_1(e_{1,1}\nu_1+e_{1,2}\nu_2)}(-1)^{\varepsilon_2(e_{2,1}\nu_1+e_{2,2}\nu_2)}\, r_2^{n_2}\\
&\times \prod_{3,\mu}^{h} \sum^{n_h} (-1)^{\nu_h(\varepsilon_1 e_{1,h}+\varepsilon_2 e_{2,h})}\, r_h^{n_h}.
\end{aligned} \]

In den fünf oben aufgestellten Gruppen, die zu $m = 24$
gehören, kann man die Perioden aus dieser Formel oder auch
leicht direct berechnen, wenn man
\[ r = e^{\frac{\pi i}{12}} = e^{\frac{\pi i}{3}}\, e^{-\frac{\pi i}{4}} \]
setzt. Man erhält dann für die fünf Grössen:
\[ \mathfrak{A}_1)\quad \eta = \frac{(1+i)\sqrt{3}}{\sqrt{2}} \]
\[ \mathfrak{A}_2)\quad \eta = \frac{\sqrt{3}-i}{\sqrt{2}} \]
%CONT>>
